% RESEARCH_DRAFT_SHA256: 78a83bf9f941da4478cc518a02ddeb48a70e20a6809e7d79d798b3f73099a32a
% Layout and source-reference wrappers expanded for the web renderer.
\documentclass[11pt]{article}
\usepackage[margin=25mm]{geometry}
\usepackage{fontspec,unicode-math}
\setmainfont{Latin Modern Roman}
\setmathfont{Latin Modern Math}
\usepackage{amsmath,amssymb,mathtools,xurl,hyperref}
\hypersetup{colorlinks=true,urlcolor=blue,linkcolor=blue}
\setcounter{secnumdepth}{0}
\setlength{\parindent}{0pt}
\setlength{\parskip}{5pt}
\setlength{\emergencystretch}{3em}
\newcommand{\sref}[2]{\hyperlink{source:#1:#2}{[S#2]}}
\newcommand{\cataloguescope}[1]{\paragraph{Source record.}#1}
\newcommand{\floor}[1]{\left\lfloor#1\right\rfloor}
\begin{document}
\section*{Research-attempt provenance}
This standalone research attempt retains the complete original definition body below,
including its historical source and preparation notes. Those notes describe the
earlier catalogue preparation; the new attempt and its verification are identified
separately after that body. No original ID, catalogue status, or source statement is changed.

\section{Prox-Regularity of Polynomial Stability Abscissas}\label{statement}

\subsection{Definitions and mathematical statement}

\paragraph{Shared notation.}Standard number systems are $\mathbb N,\mathbb Z,\mathbb Q,\mathbb R,\mathbb C$. Specialized notation and index ranges follow the cited source and the definitions below.


The stability abscissa of $p$ is $\alpha(p)=\max\{\operatorname{Re}z:p(z)=0\}$. Stability concerns roots in the stated left half-plane; a nearest stable polynomial minimizes the chosen coefficient-space distance.
\paragraph{Record scope and qualifications.}
The catalogue title is retained as metadata. Named certificates, contracts, bounds or frameworks in that title are not asserted as results of the original source extract.
The coefficient norm, meaning of near, treatment of roots on the imaginary axis, and exact prox-regularity convention are not specified. The source\textquotesingle s asserted equivalence is retained rather than independently established.

Source-record limitation: The distance norm and prox-regularity convention defining the claimed equivalence are absent from the supplied formulation.
\paragraph{Statement recorded in the supplied source.}
Does every degree-$k$ polynomial $p(s)$ near $s^k$ have a unique nearest stable polynomial? Equivalently, is the stability abscissa $\alpha$ prox-regular there? \href{https://doi.org/10.4171/owr/2005/02}{[S2]}
\paragraph{Source status.}
The supplied catalogue labels this source record open; that metadata was checked against the supplied snapshot on 2026-10-11. The label is preserved for provenance and is not an independent verification of present open status or mathematical correctness. The catalogue literature assessment is dated 2026-08-21.
\subsection{Short English statement}
Decide whether nearby polynomials have unique closest stable approximations and whether the stated regularity property holds.
\subsection{Sources}
\begin{itemize}
\item[S1] Ulam AI and the UnsolvedMath Contributors, UnsolvedMath: A Curated Collection of Open Mathematics Problems, catalogue entry. Catalogue curation and metadata: CC BY 4.0. The source extract is retained; definitions, English summaries and qualifications are editorial additions.
\url{https://huggingface.co/datasets/ulamai/UnsolvedMath}.
\item[S2] Optimization and Applications (2005), Oberwolfach Reports, DOI 10.4171/owr/2005/02, pp. 69--133
\url{https://doi.org/10.4171/owr/2005/02}.
\item[S3] Burke--Lewis--Overton
\url{https://sites.math.washington.edu/~burke/papers/reprints/53-va-polys-GL.pdf}.
\item[S4] primary OWR paper text
\url{https://ems.press/content/serial-article-files/45976?nt=1}.
\item Source-specific Oberwolfach Reports licences, including CC BY-SA 4.0 where applicable, are retained; see THIRD\_PARTY\_NOTICES.md.
\end{itemize}
\paragraph{Preparation and review.}
The original source and UnsolvedMath contributors retain attribution for the source extract. Codex (GPT-6; exact serving revision not exposed) assisted with definitions, English summaries, scope notes and formatting. No new human mathematical authorship or independently verified solution is asserted. Catalogue identity checked against the complete supplied registry and current repository. Distinct upstream records, including repeated questions and broad or damaged extracts, are retained for complete coverage of the source CS category. A repeated record is not a new independent conjecture. Complete mathematical-equivalence and current-literature review has not been performed. Preparation and final review follow contribution prompts C and D, with source gaps disclosed. Structural and rendering checks do not establish a new theorem or open status.
\label{end:statement}

\section{Literature and context check}
The primary report, printed page 106, discusses nearest projection onto an epigraph and calls the nearest-stable-polynomial question an essential ingredient. It does not state the catalogue's explicit equivalence as a proved theorem. That stronger wording is preserved above, but is not used below. To make a concrete calculation, this attempt restricts to real monic quadratics, Euclidean distance on their two nonleading coefficients, and closed left-half-plane stability. These choices are auxiliary scope, not a repair of the full degree-$k$ problem.
\section{Research attempt: Closed quadratic stability and the limits of a root-sensitivity argument}

\paragraph{An exact projection calculation}
For $p(s)=s^2+as+b$ with real $a,b$, both roots have real part at most zero if and only if $a\geq0$ and $b\geq0$. Indeed, real roots with these sum and product conditions have sum $-a\leq0$ and product $b\geq0$, so they are both nonpositive (including a possible zero). Nonreal roots are a conjugate pair with common real part $-a/2$ and product equal to their squared modulus. These observations also prove necessity, considering the real and nonreal cases separately.

Consequently the closed stable coefficient set is exactly the quadrant $[0,\infty)^2$. The nearest stable polynomial in the stated Euclidean norm has coefficients
\[
 (a_*,b_*)=(\max(a,0),\max(b,0)).
\]
To prove uniqueness directly, the squared distance splits as $(u-a)^2+(v-b)^2$ with independent constraints $u,v\geq0$. Each term is uniquely minimized at the corresponding truncated coordinate: if its center is nonnegative it is the unconstrained minimum, and otherwise the square strictly increases on the nonnegative half-line. Thus this degree-two closed-real case has unique projection everywhere, in particular near $s^2$.

\paragraph{A failed attempt to infer a negative answer from root sensitivity}
The root formula gives the exact abscissa
\[
 \alpha(a,b)=-a/2+\sqrt{\max(a^2/4-b,0)}.
\]
Along $p_\varepsilon(s)=s^2-\varepsilon$, $\varepsilon>0$, coefficient distance from $s^2$ is $\varepsilon$, whereas abscissa changes from zero to $\sqrt\varepsilon$. Hence the abscissa is not locally Lipschitz at the origin in this coefficient space. This does not refute uniqueness of nearest closed-stable projection: the preceding explicit calculation proves uniqueness for the very same family. Nor has any theorem equating lack of local Lipschitz continuity with failure of the requested prox-regularity been supplied. The singular root derivative therefore cannot by itself settle the report's question.

\paragraph{Why the stability convention matters}
For the same real quadratics, strictly negative real parts are equivalent to $a>0,b>0$, by the same root argument with strict inequalities. This set is open. At $p=s^2$ the strictly stable sequence $s^2+\delta s+\delta$, $\delta>0$, approaches coefficient distance zero, but no strictly stable polynomial attains zero distance. More generally, from any strictly stable candidate for an exterior point one can move a sufficiently small positive distance toward that point and remain in the open quadrant, while strictly reducing the distance. Thus there is no nearest strictly stable polynomial for an exterior point. This is an ambiguity check, not a counterexample to an intended formulation with the closed stable set.

\section{Verification and adversarial audit}
The real-root case includes double roots, $a=0$, $b=0$, and the origin; the conjugate-root case includes roots on the imaginary axis. The projection concerns the closed stable coefficient set, whereas the source's prox-regularity discussion concerns an epigraph in one higher dimension. These sets are not identified. The exact square-root example exposes failure of one inference while retaining the directly proved projection result.
\paragraph{Audit corrections.} Distinguished closed and strict stability, stable-set projection and epigraph projection, and non-Lipschitz root behavior from a prox-regularity conclusion.
\paragraph{External computation.} None; symbolic arguments and boundary checks only.
\section{Scope and remaining gaps}
General degree, complex coefficients, other norms, and the source's precise prox-regularity convention are not settled. The purported equivalence in the retained catalogue is not adopted without proof. The closed-real quadratic case and the open-set obstruction are known elementary material rederived to isolate the missing scope.
\paragraph{Integration review.} Codex (GPT-6; exact serving revision not exposed) reviewed the written arguments and expanded layout and source-reference wrappers. Repository metadata and the production web renderer were checked using Python 3.13 and JavaScript with MathJax 3.2.2. This model-assisted review is not independent human proof verification; no PDF compilation or formal proof checking is asserted.

\section{Completion Estimate}
\noindent\textbf{COMPLETION: 0\%}
The calculations below the preserved source statement investigate a concrete route,
but establish no substantiated new progress toward the entire original target.
Known special cases and failed sufficient conditions are not credited as a solution.
\section{References}
The original source references and catalogue attribution are preserved in the
Sources section above. The following URLs were actually consulted for this attempt;
retention of an additional original reference does not claim fresh consultation of it.
\begin{itemize}
\item \url{https://ems.press/content/serial-article-files/45976} (consulted 2026-10-11). 
\end{itemize}
\section{Model and contributor disclosure}
OpenAI gpt-6-astra, requested reasoning configuration ultra, generated this research
attempt on 2026-10-11. The exact serving revision is not exposed. Codex GPT-6
performed orchestration and repository-format review; its exact serving revision
is also not exposed. The mathematical route, boundary checks and final adversarial
self-audit followed contribution prompts A then B in
\texttt{docs/contribution\_prompts.md}, after reading \texttt{CONTRIBUTING.md}.
Python 3.14.2 assembled the standalone documents; any mathematical computation is
specified in Verification. No theorem prover was used. No human mathematical
authorship, substantial human proof correction, or independent human proof checking
is asserted. Model-assisted review, source retrieval, and structural checks do not
establish correctness or novelty. Full private reasoning transcripts are unavailable
as reproducibility artifacts; the complete arguments, inputs and outcomes actually
used are stated here or in the companion batch audit.
\section{Final status}
Unresolved as a new research contribution. This is a concrete route investigation,
not a claim of new complete resolution or a current-status certification. A known
answer reproduced for context does not become a new contribution. The catalogue's
source-status label is preserved; this attempt status does not assert that a
source-labelled solved question is open.
\end{document}
